\subsection{Quantitative Comparison on KITTI Benchmark}
\label{sec:benchmark_comparison}

\begin{table*}[t]
\centering\small
\caption{Multi-seed statistical robustness and variance comparison between baseline MobilePIXOR and the proposed MobilePIXORNeXt on the KITTI validation benchmark across three random seeds ($42, 43, 44$) under a 100-epoch training schedule. Detections are evaluated under the unified BEV evaluation protocol with 40-point interpolation ($R_{40}$) on NVIDIA L4. Performance is reported across individual seeds alongside the sample mean and standard deviation ($\text{Mean} \pm \text{Std}$, \%). All three baseline runs employ the canonical Binary35 input representation, ensuring an unbiased, standardized multi-seed baseline. The best mean and lowest variance values are highlighted in \textbf{bold}.}
\label{tab:seed_robustness}
\setlength{\tabcolsep}{4.5pt}
\begin{tabular}{@{}lrrrrrrrrrrr@{}}
\toprule
& \multicolumn{5}{c}{\textbf{Baseline MobilePIXOR~\cite{nguyen2024indoor}}} & \multicolumn{5}{c}{\textbf{MobilePIXORNeXt (Ours)}} & \\
\cmidrule(lr){2-6} \cmidrule(lr){7-11}
\textbf{Category / Metric} & \textbf{S42} & \textbf{S43} & \textbf{S44} & \textbf{Mean} & \textbf{Std} & \textbf{S42} & \textbf{S43} & \textbf{S44} & \textbf{Mean} & \textbf{Std} & \textbf{Gain ($\Delta$)} \\
\midrule
Car AP$_{\mathrm{M}}$ (\%) & 75.52 & 66.75 & 71.83 & 71.37 & $\pm 4.40$ & 86.90 & 86.33 & 85.42 & \textbf{86.22} & $\mathbf{\pm 0.75}$ & \textit{+14.85} \\
Pedestrian AP$_{\mathrm{M}}$ (\%) & 28.16 & 28.54 & 26.56 & 27.75 & $\pm 1.05$ & 60.78 & 61.94 & 60.06 & \textbf{60.93} & $\mathbf{\pm 0.95}$ & \textit{\textbf{+33.17}} \\
Cyclist AP$_{\mathrm{M}}$ (\%) & 20.53 & 26.60 & 18.73 & 21.95 & $\pm 4.12$ & 66.81 & 65.96 & 64.26 & \textbf{65.68} & $\mathbf{\pm 1.30}$ & \textit{\textbf{+43.72}} \\
\midrule
\textbf{Moderate mAP (\%)} & 41.40 & 40.63 & 39.04 & 40.36 & $\pm 1.20$ & 71.50 & 71.41 & 69.91 & \textbf{70.94} & $\mathbf{\pm 0.89}$ & \textit{\textbf{+30.58}} \\
\textbf{Overall mAP$_{9}$ (\%)} & 42.27 & 41.77 & 41.24 & 41.76 & $\pm 0.52$ & 73.45 & 73.35 & 72.22 & \textbf{73.01} & $\mathbf{\pm 0.68}$ & \textit{\textbf{+31.25}} \\
\bottomrule
\end{tabular}
\par\vspace{2pt}
{\footnotesize \textit{Note.} All three baseline MobilePIXOR configurations are standardized to use the canonical Binary35 (\texttt{binary\_slices}) BEV encoding and trained with the identical AdamW optimizer on NVIDIA L4, correcting earlier input representation mismatches across seeds. MobilePIXORNeXt demonstrates superior stability ($\pm 0.89\%$ vs. $\pm 1.20\%$ Moderate mAP) and delivers substantial gains of $+30.58\,\mathrm{pp}$ Moderate mAP and $+31.25\,\mathrm{pp}$ $\mathrm{mAP}_9$.}
\end{table*}


\paragraph{Detection Quality and Parameter Efficiency.}
Table~\ref{tab:comparison} benchmarks MobilePIXORNeXt against six established detectors on the KITTI validation split using the Unified KITTI BEV Evaluation Protocol described in Section~\ref{sec:dataset_and_metrics}. Detections are scored with identical ground-truth loading, spatial ROI filtering ($0\text{--}70.4\,\text{m}$ forward range), class-neighbour ignore policies, and 40-point interpolation ($R_{40}$).

\begin{figure*}[t]
\centering
\paperfigure[width=\textwidth]{pareto_tradeoff.pdf}{48mm}{L4 accuracy, parameter count, and pipeline latency.}
\caption{Accuracy--efficiency comparison on KITTI validation: Moderate BEV mAP ($R_{40}$) versus (a) model parameter count and (b) L4 end-to-end pipeline latency. Dashed lines connect nondominated configurations among the reported models.}
\label{fig:pareto_tradeoff}
\end{figure*}

\textit{1. Accuracy and Parameter Efficiency.}

MobilePIXORNeXt contains $753{,}569$ deployable parameters ($0.754\,\text{M}$) and 12 training-only loss-balancing logits. The seed-42 epoch-100 checkpoint reaches $71.50\%$ Moderate mAP and $73.45\%$ $\mathrm{mAP}_9$. It exceeds PointPillars ($70.27\%$) by $1.22\,\text{pp}$ with $84.41\%$ fewer parameters, and trails SECOND ($72.18\%$) by $0.68\,\text{pp}$ with $85.85\%$ fewer parameters. Figure~\ref{fig:pareto_tradeoff} illustrates both parameter efficiency and end-to-end pipeline latency, demonstrating that MobilePIXORNeXt lies on the Pareto frontier in both regimes and achieves faster end-to-end execution than SECOND. Figure~\ref{fig:pr_curves} details the full precision--recall trajectories across all nine class--difficulty evaluation regimes.

\begin{figure*}[!t]
\centering
\paperfigure[width=\textwidth]{pr_all_difficulties_sota.pdf}{65mm}{Precision-Recall benchmark curves across classes and difficulties on the KITTI validation set.}
\caption{Precision--recall curves on the KITTI validation benchmark ($R_{40}$). Columns display Car ($\text{IoU} \ge 0.70$), Pedestrian ($\text{IoU} \ge 0.50$), and Cyclist ($\text{IoU} \ge 0.50$); rows correspond to Easy, Moderate, and Hard difficulties. MobilePIXORNeXt (Ours) exhibits consistently high precision and broad recall coverage across all nine evaluation regimes, substantially outperforming prior compact detectors especially on vulnerable road users.}
\label{fig:pr_curves}
\end{figure*}

\textit{2. Performance on Pedestrians.}

MobilePIXORNeXt reaches $60.78\%$ Pedestrian Moderate AP versus $28.16\%$ for the baseline MobilePIXOR checkpoint, representing a substantial gain of $+32.62\,\text{pp}$. Furthermore, its score surpasses established 3D anchor and pillar architectures, exceeding PointPillars ($55.83\%$) and SECOND ($56.47\%$) by $+4.95\,\text{pp}$ and $+4.31\,\text{pp}$, respectively, trailing only the heavyweight two-stage Part-A$^2$-Anchor ($63.56\%$, which requires $63.8\,\text{M}$ parameters).

\textit{3. Multi-Class Detection and Baseline Comparison.}

MobilePIXORNeXt records $86.90\%$ Car and $66.81\%$ Cyclist Moderate AP. Its $73.45\%$ overall nine-cell mean ($\text{mAP}_9$) exceeds PointPillars ($71.80\%$) and SECOND ($73.16\%$), confirming robust detection coverage across diverse object categories and difficulty levels. Compared with baseline MobilePIXOR ($41.40\%$ Moderate mAP and $42.27\%$ $\text{mAP}_9$), MobilePIXORNeXt achieves seed-42 improvements of $+30.10\,\text{pp}$ and $+31.18\,\text{pp}$, demonstrating the efficacy of our integrated architecture.

\textit{4. Multi-Seed Statistical Robustness and Variance Analysis.}

Table~\ref{tab:seed_robustness} summarizes the multi-seed evaluation across random seeds ($42, 43, 44$) for both baseline MobilePIXOR and the proposed MobilePIXORNeXt. Baseline MobilePIXOR runs are standardized to the canonical Binary35 representation (\texttt{binary\_slices}), correcting earlier input representation discrepancies across seeds. Under this standardized baseline, MobilePIXOR achieves a mean Moderate mAP of $40.36\pm1.20\%$ ($71.37\pm4.40\%$ Car, $27.75\pm1.05\%$ Pedestrian, $21.95\pm4.12\%$ Cyclist) and $\text{mAP}_9$ of $41.76\pm0.52\%$. In contrast, MobilePIXORNeXt demonstrates superior accuracy and remarkable stability, achieving a mean Moderate mAP of $70.94\pm0.89\%$ ($86.22\pm0.75\%$ Car, $60.93\pm0.95\%$ Pedestrian, $65.68\pm1.30\%$ Cyclist) and $\text{mAP}_9$ of $73.01\pm0.68\%$. 

The resulting net performance gains are substantial: $+14.85\,\text{pp}$ on Car, $+33.17\,\text{pp}$ on Pedestrian, $+43.72\,\text{pp}$ on Cyclist, $+30.58\,\text{pp}$ Moderate mAP, and $+31.25\,\text{pp}$ overall $\text{mAP}_9$. Crucially, standardizing the baseline input representation on Binary35 eliminates the previous high-variance anomaly ($\pm 5.06\%$ under mixed representations reduced to $\pm 1.20\%$). MobilePIXORNeXt further suppresses variance across seeds ($\pm 0.89\%$ vs. $\pm 1.20\%$ Moderate mAP, and $\pm 0.75\%$ vs. $\pm 4.40\%$ on Car AP$_{\mathrm{M}}$), confirming that the architectural combination of Strata-BEV14, inverted bottleneck residual blocks, Block-SimAM, and bounded multi-task optimization (T-SBUW) produces robust, reproducible convergence regardless of random weight initialization.

\subsection{Inference Latency and Real-Time Complexity}
\label{sec:latency_analysis}
\begin{table*}[t]
\centering\small
\caption{Offline pipeline latency and throughput on NVIDIA L4, batch size 1, FP32. Accuracy uses all 3,769 validation frames; latency uses 3,759 frames after 10 warmup frames. Best and second-best values among compact models are bold and underlined.}
\label{tab:efficiency}
\setlength{\tabcolsep}{3pt}
\begin{tabular}{@{}lrrrrrrrr@{}}
\toprule
\textbf{Method} & \textbf{\shortstack{Params\\(M) $\downarrow$}} & \textbf{\shortstack{Model\\(ms) $\downarrow$}} & \textbf{\shortstack{Model\\FPS $\uparrow$}} & \textbf{\shortstack{E2E\\(ms) $\downarrow$}} & \textbf{\shortstack{E2E\\FPS $\uparrow$}} & \textbf{\shortstack{P50\\(ms) $\downarrow$}} & \textbf{\shortstack{P95\\(ms) $\downarrow$}} & \textbf{\shortstack{P99\\(ms) $\downarrow$}} \\
\midrule
MobilePIXOR~\cite{nguyen2024indoor} & \textbf{0.598} & \underline{11.04} & \underline{90.57} & 88.02 & 11.36 & 86.04 & 116.13 & 130.45 \\
PointPillars~\cite{lang2019pointpillars} & 4.835 & 12.04 & 83.02 & \textbf{27.48} & \textbf{36.40} & \textbf{27.03} & \textbf{31.11} & \textbf{33.10} \\
SECOND~\cite{yan2018second} & 5.326 & 26.39 & 37.89 & 40.83 & 24.49 & 40.91 & \underline{44.31} & \underline{48.43} \\
\textbf{MobilePIXORNeXt (Ours)} & \underline{0.754} & \textbf{8.67} & \textbf{115.32} & \underline{37.29} & \underline{26.82} & \underline{36.05} & 46.98 & 51.32 \\
\midrule
\multicolumn{9}{@{}l}{\textit{Larger 3D Reference Detectors}} \\
PointRCNN~\cite{shi2019pointrcnn} & 4.039 & 79.48 & 12.58 & 90.86 & 11.01 & 90.90 & 95.04 & 97.00 \\
Part-A$^2$-Anchor~\cite{shi2020part} & 63.815 & 61.81 & 16.18 & 76.81 & 13.02 & 76.65 & 81.77 & 86.10 \\
PV-RCNN~\cite{shi2020pvrcnn} & 13.116 & 101.36 & 9.87 & 114.08 & 8.77 & 113.89 & 119.91 & 122.41 \\
\bottomrule
\end{tabular}
\par\vspace{2pt}
{\footnotesize \textit{Note.} End-to-end (E2E) pipeline latency encompasses point-cloud file loading, CPU input rasterization, H2D memory transfer, network execution, and rotated NMS decoding. All models are benchmarked under sequential batch-1 execution on a dedicated NVIDIA L4 GPU in FP32 precision. FPS is calculated as $1000/\text{mean latency}$.}
\end{table*}


\paragraph{Pipeline Latency vs. Model-Stage Throughput.}
All current benchmark reports use NVIDIA L4, batch size 1, and FP32 inference. Accuracy covers 3,769 frames; timing excludes the first 10 warmup frames. MobilePIXORNeXt's model stage takes $8.67\,\text{ms}$ ($115.32\,\text{FPS}$), while its complete sequential pipeline takes $37.29\,\text{ms}$ ($26.82\,\text{FPS}$). The network-only value is a stage within the pipeline, not an independently measured standalone operating point. External-detector model stages include architecture-specific internal decoding and proposal NMS, whereas final post-processing is measured separately.

\textit{1. End-to-End Pipeline Comparison.}

MobilePIXOR requires $88.02\,\text{ms}$ ($11.36\,\text{FPS}$); in comparison, MobilePIXORNeXt cuts mean pipeline latency by $57.63\%$ to $37.29\,\text{ms}$ ($26.82\,\text{FPS}$), surpassing SECOND ($40.83\,\text{ms}$, $24.49\,\text{FPS}$) to achieve the second-fastest end-to-end throughput among compact detectors behind PointPillars ($27.48\,\text{ms}$, $36.40\,\text{FPS}$). MobilePIXORNeXt exhibits stable real-time execution with a median latency of $36.05\,\text{ms}$, while tail percentiles ($46.98\,\text{ms}$ P95 and $51.32\,\text{ms}$ P99) reflect CPU rasterization on densely populated point clouds, highlighting pre-processing parallelization as a viable avenue for future embedded acceleration.

\textit{2. Pipeline Components and Transfer Cost.}

For MobilePIXORNeXt, the four stage means sum to $37.29\,\text{ms}$:
\begin{itemize}
    \item \textbf{Input preparation ($11.64\,\text{ms}$)} includes point-cloud file loading and Strata-BEV14 rasterization with the Numba CPU backend. This is not an encoder-only timing. PointPillars and SECOND record $12.67\,\text{ms}$ and $11.92\,\text{ms}$ for their native input preparation.
    \item \textbf{H2D transfer ($6.67\,\text{ms}$)} moves the $31.54\,\text{MB}$ FP32 input. Binary35 requires $78.85\,\text{MB}$, so Strata-BEV14 reduces the tensor payload by $60.0\%$. MobilePIXOR's measured transfer is $41.75\,\text{ms}$, but implementation and memory handling also differ; the timing ratio is not a pure channel-count effect.
    \item \textbf{Model stage ($8.67\,\text{ms}$)} is lower than MobilePIXOR ($11.04\,\text{ms}$), PointPillars ($12.04\,\text{ms}$), and SECOND ($26.39\,\text{ms}$), subject to their different internal stage boundaries.
    \item \textbf{Decoding and rotated NMS ($10.31\,\text{ms}$)} exceed the external PointPillars final-post-processing stage ($1.40\,\text{ms}$). These implementations partition decoding differently. Input transfer and post-processing are therefore important targets for pipeline optimization.
\end{itemize}

\subsection{Distance-Stratified Detection Performance}
\label{sec:distance_stratified}

\paragraph{Range-Dependent Accuracy on Vulnerable Road Users.}
Table~\ref{tab:distance_ap} evaluates detection performance across three radial distance intervals: Near ($0\text{--}30\,\text{m}$), Intermediate ($30\text{--}50\,\text{m}$), and Far ($50\text{--}70.4\,\text{m}$) on the KITTI validation benchmark using Moderate BEV AP $R_{40}$.

\begin{table}[!htbp]
\centering\footnotesize
\caption{Distance-stratified detection performance (Moderate AP, \%) on Vulnerable Road Users (Pedestrian and Cyclist) across Near ($0\text{--}30\,\text{m}$), Intermediate ($30\text{--}50\,\text{m}$), and Far ($50\text{--}70.4\,\text{m}$) range intervals on the KITTI validation benchmark using BEV AP $R_{40}$. The supplied seed-42 L4 checkpoints are used. Best results are in \textbf{bold}.}
\label{tab:distance_ap}
\setlength{\tabcolsep}{1.5pt}
\begin{tabular*}{\columnwidth}{@{\extracolsep{\fill}}llrrr@{}}
\toprule
\textbf{Category} & \textbf{Model} & \textbf{\shortstack{0--30\,m}} & \textbf{\shortstack{30--50\,m}} & \textbf{\shortstack{50--70.4\,m}} \\
\midrule
Pedestrian & MobilePIXOR~\cite{nguyen2024indoor} & 30.45 & 3.75 & \textbf{0.99} \\
& \textbf{Ours} & \textbf{65.26} & \textbf{11.56} & 0.80 \\
& \textit{Diff.} & \textit{+34.81} & \textit{+7.81} & \textit{-0.19} \\
\midrule
Cyclist & MobilePIXOR~\cite{nguyen2024indoor} & 25.55 & 12.61 & 0.56 \\
& \textbf{Ours} & \textbf{76.32} & \textbf{44.33} & \textbf{15.45} \\
& \textit{Diff.} & \textit{+50.76} & \textit{+31.72} & \textit{+14.89} \\
\bottomrule
\end{tabular*}
\end{table}


\textit{1. Near Range ($0\text{--}30\,\text{m}$).}

MobilePIXORNeXt records $65.26\%$ Pedestrian AP versus $30.45\%$ for MobilePIXOR ($+34.81\,\text{pp}$), and $76.32\%$ Cyclist AP versus $25.55\%$ ($+50.76\,\text{pp}$). Differences are computed from unrounded report values.

\textit{2. Intermediate and Long Ranges ($30\text{--}70.4\,\text{m}$).}

At $30\text{--}50\,\text{m}$, Pedestrian AP is $11.56\%$ versus $3.75\%$, and Cyclist AP is $44.33\%$ versus $12.61\%$. At $50\text{--}70.4\,\text{m}$, Cyclist AP is $15.45\%$ versus $0.56\%$ ($+14.89\,\text{pp}$). Far-range Pedestrian AP decreases from $0.99\%$ to $0.80\%$ ($-0.19\,\text{pp}$). Far-range ground-truth counts are only 13 Pedestrians and 17 Cyclists, so those AP values should be interpreted with this limited support in mind.

\begin{table*}[t]
\centering\small
\caption{Architectural progression roadmap on KITTI validation using Moderate BEV AP ($R_{40}$). Panel (a) outlines the progressive single-factor integration of core proposed modules from baseline MobilePIXOR ($A_0$) to the final MobilePIXORNeXt configuration ($A_8$). Panel (b) reports cumulative detection accuracy ($\text{mAP}_{\mathrm{M}}$, $\text{mAP}_9$) and model-stage inference latency and throughput (ms, FPS) measured on NVIDIA L4. Best values are in \textbf{bold}, and second-best are \underline{underlined}.}
\label{tab:components}
\textbf{(a) Core architectural milestones and component progression}\par\smallskip
\begingroup\footnotesize\setlength{\tabcolsep}{1.5pt}
\begin{tabular*}{\textwidth}{@{\extracolsep{\fill}}lllllll@{}}
\toprule
\textbf{Step} & \textbf{Backbone} & \textbf{BEV Input} & \textbf{Head} & \textbf{Augmentation} & \textbf{Neck} & \textbf{Loss / Attention} \\
\midrule
$A_0$ & MobilePIXOR & Binary35 & Single Head & Standard Aug & FPN & Baseline Smooth-L1 \\
$A_1$ & MobilePIXORNeXt & Binary35 & Single Head & Standard Aug & FPN & Baseline Smooth-L1 \\
$A_2$ & MobilePIXORNeXt & Strata-BEV14 & Single Head & Standard Aug & FPN & Baseline Smooth-L1 \\
$A_3$ & MobilePIXORNeXt & Strata-BEV14 & Decoupled Grouped & Standard Aug & FPN & Baseline Smooth-L1 \\
$A_4$ & MobilePIXORNeXt & Strata-BEV14 & Decoupled Grouped & Standard Aug & FPN & Baseline Smooth-L1 \\
$A_5$ & MobilePIXORNeXt & Strata-BEV14 & Decoupled Grouped & PC-GTA & FPN & Baseline Smooth-L1 \\
$A_6$ & MobilePIXORNeXt & Strata-BEV14 & Decoupled Grouped & PC-GTA & SG-FPN & Baseline Smooth-L1 \\
$A_7$ & MobilePIXORNeXt & Strata-BEV14 & Decoupled Grouped & PC-GTA & SG-FPN & OGA + IQA Head \\
$A_8$ & MobilePIXORNeXt & Strata-BEV14 & Decoupled Grouped & PC-GTA & SG-FPN & OGA + IQA + Block-SimAM \\
\bottomrule
\end{tabular*}\endgroup
\par\medskip
\textbf{(b) Quantitative detection accuracy and model-stage inference speed}\par\smallskip
\begingroup\footnotesize\setlength{\tabcolsep}{3.0pt}
\begin{tabular*}{\textwidth}{@{\extracolsep{\fill}}l>{\raggedright\arraybackslash}p{0.45\textwidth}rrrr@{}}
\toprule
\textbf{Step} & \textbf{Investigated Milestone / Architectural Focus} & \textbf{\shortstack{mAP$_{\mathrm{M}}$\\(\%) $\uparrow$}} & \textbf{\shortstack{mAP$_{9}$\\(\%) $\uparrow$}} & \textbf{\shortstack{Model Stage\\(ms) $\downarrow$}} & \textbf{\shortstack{Model FPS\\$\uparrow$}} \\
\midrule
$A_0$ & Baseline MobilePIXOR starting checkpoint & 41.40 & 42.27 & 11.04 & 90.57 \\
$A_1$ & + ConvNeXt-style backbone ($7\times 7$ DW, SiLU, Inverted Bottleneck) & 47.60 & 49.08 & \textbf{5.88} & \textbf{169.97} \\
$A_2$ & + Strata-BEV14 statistical representation (subcell offsets, height moments) & 50.14 & 52.00 & \underline{5.92} & \underline{169.01} \\
$A_3$ & + Decoupled GroupedHeader (Car vs. Ped/Cyc independent towers) & 51.30 & 53.96 & 6.92 & 144.58 \\
$A_4$ & + Expanded prediction head capacity (wider representation channels) & 60.22 & 61.11 & 6.70 & 149.15 \\
$A_5$ & + Physics-Consistent Ground-Truth Augmentation (PC-GTA) & 68.12 & 70.05 & 6.72 & 148.82 \\
$A_6$ & + Scale-Gated Feature Pyramid Network (SG-FPN) & 68.46 & 70.28 & 6.94 & 144.16 \\
$A_7$ & + Oriented Geometric Alignment (OGA) and IQA quality ranking & \underline{70.58} & \underline{72.66} & 8.19 & 122.16 \\
$A_8$ & + Parameter-Free Block-SimAM Energy Attention inside all blocks (Final) & \textbf{71.50} & \textbf{73.45} & 8.67 & 115.32 \\
\bottomrule
\end{tabular*}\endgroup
\par\vspace{2pt}
{\footnotesize \textit{Note.} Step identifiers $A_0$--$A_8$ illustrate the progressive single-factor architectural evolution from baseline MobilePIXOR ($A_0$) to the proposed MobilePIXORNeXt ($A_8$). Each step isolates a single architectural modification. Controlled ablations isolating individual loss objectives and attention mechanisms under fixed experimental budgets are reported separately in Table~\ref{tab:loss} (Loss Matrix) and Table~\ref{tab:attention} (Attention Placement). Timings are profiled on a dedicated NVIDIA L4 GPU under batch size 1 in FP32 precision.}
\end{table*}


\subsection{Ablation Studies}
\label{sec:ablations}

\paragraph{Systematic Single-Factor Architectural Progression.}
Table~\ref{tab:components} outlines the cumulative progression roadmap of architectural refinements bridging the baseline MobilePIXOR ($A_0$) to the proposed MobilePIXORNeXt configuration ($A_8$). Each step isolates a single architectural modification, demonstrating its empirical contribution:

\paragraph{1. Backbone Architecture ($A_0 \to A_1$).}
Replacing MobilePIXOR's legacy residual blocks with $7\times 7$ depthwise inverted bottleneck blocks (\texttt{MobilePIXORNeXt-S}, $e=2.5$, SiLU, LayerScale) yields a marked $+6.20\,\text{pp}$ increase in Moderate mAP ($41.40\% \to 47.60\%$) while almost halving forward execution latency ($11.04\,\text{ms} \to 5.88\,\text{ms}$, $+87.7\%$ model throughput).

\paragraph{2. Statistical BEV Representation ($A_1 \to A_2$).}
\begin{table*}[t]
\centering\small
\caption{Ablation of Bird's-Eye View (BEV) input representation paradigms on the KITTI validation split using the local BEV evaluator with $R_{40}$ interpolation (Section~\ref{sec:dataset_and_metrics}). The listed models are benchmarked under the standardized baseline configuration (MobilePIXORNeXt backbone, SumFPN neck, single head, standard augmentation, baseline loss, seed 42). Best results across all configurations are in \textbf{bold}, and second-best are \underline{underlined}.}
\label{tab:encoding}
\setlength{\tabcolsep}{2.5pt}
\begin{tabular*}{\textwidth}{@{\extracolsep{\fill}}lcrrrrrrr@{}}
\toprule
\textbf{Representation} & \textbf{$C_{\mathrm{in}}$} & \textbf{\shortstack{H2D Payload\\(MB) $\downarrow$}} & \textbf{\shortstack{Car AP$_{\mathrm{M}}$\\(\%) $\uparrow$}} & \textbf{\shortstack{Pedestrian\\AP$_{\mathrm{M}}$ (\%) $\uparrow$}} & \textbf{\shortstack{Cyclist\\AP$_{\mathrm{M}}$ (\%) $\uparrow$}} & \textbf{\shortstack{mAP$_{\mathrm{M}}$\\(\%) $\uparrow$}} & \textbf{\shortstack{mAP$_{9}$\\(\%) $\uparrow$}} & \textbf{\shortstack{Standalone\\FPS $\uparrow$}} \\
\midrule
TriBand-BEV~\cite{khoshkdahan2026triband} & 3 & \textbf{6.76} & 74.10 & 38.05 & 21.12 & 44.42 & 46.55 & \textbf{188.43} \\
Strata-BEV8  & 8 & \underline{18.02} & \underline{78.35} & 38.48 & \textbf{25.29} & 47.37 & 48.95 & 165.02 \\
Binary35~\cite{yang2018pixor} & 35 & 78.85 & 75.24 & \underline{44.36} & 23.20 & \underline{47.60} & \underline{49.08} & \underline{169.97} \\
\textbf{Strata-BEV14 (Ours)} & \textbf{14} & 31.54 & \textbf{79.97} & \textbf{45.36} & \underline{25.10} & \textbf{50.14} & \textbf{52.00} & 169.01 \\
\bottomrule
\end{tabular*}
\par\vspace{2pt}
{\footnotesize \textit{Note.} All configurations are evaluated under the baseline architecture (MobilePIXORNeXt backbone, SumFPN neck, single head, standard augmentation, baseline loss, seed 42). Payloads assume an $800\times704$ FP32 tensor grid. Strata-BEV14 achieves the highest detection accuracy while reducing host-to-device memory transfer payload by $60.0\%$ relative to Binary35 ($31.54\,\mathrm{MB}$ vs. $78.85\,\mathrm{MB}$).}
\end{table*}

Transitioning from standard binary occupancy grids (Binary35) to the parameter-free Strata-BEV14 tensor advances Moderate mAP from $47.60\%$ to $50.14\%$ ($+2.54\,\text{pp}$) while slashing host-to-device memory transfer payload by $60.0\%$ ($31.54\,\text{MB}$ vs. $78.85\,\text{MB}$). Continuous elevation moments and subcell centroid offsets preserve fine vertical geometry without learned point encoders. Table~\ref{tab:encoding} details the isolated encoding comparisons across intermediate variants.

\paragraph{3. Decoupled Prediction Heads ($A_2 \to A_3$).}
Decoupling prediction heads into independent Car and Pedestrian/Cyclist towers (GroupedHeader) resolves multi-task gradient interference between disparate target categories, improving Moderate mAP to $51.30\%$ ($+1.16\,\text{pp}$).

\paragraph{4. Expanded Head Capacity ($A_3 \to A_4$).}
Expanding intermediate representation width within the grouped detection towers substantially increases geometric feature capacity, lifting Moderate mAP from $51.30\%$ to $60.22\%$ ($+8.92\,\text{pp}$).

\paragraph{5. Physics-Consistent Ground-Truth Augmentation ($A_4 \to A_5$).}
Incorporating Physics-Consistent Ground-Truth Augmentation (PC-GTA) enforces local terrain elevation flushness, static obstacle clearance, and range-adaptive point downsampling, elevating Moderate mAP from $60.22\%$ to $68.12\%$ ($+7.90\,\text{pp}$) by eliminating physically implausible synthesis artifacts during training.

\paragraph{6. Scale-Gated Feature Pyramid Network ($A_5 \to A_6$).}
Introducing the bilinear depthwise Scale-Gated Feature Pyramid Network (SG-FPN) effectively mitigates transposed-convolution checkerboard artifacts and dynamically balances multi-scale spatial details, further advancing Moderate mAP to $68.46\%$ ($+0.34\,\text{pp}$).

\paragraph{7. Oriented Geometric Alignment and Quality Supervision ($A_6 \to A_7$).}
Replacing standard Smooth-L1 regression with Oriented Geometric Alignment (OGA) and localization-quality scoring (IQA) drives Moderate mAP to $70.58\%$ ($+2.12\,\text{pp}$), demonstrating the necessity of joint continuous geometric overlap feedback and quality calibration.

\paragraph{8. Parameter-Free Energy Attention ($A_7 \to A_8$).}
Embedding parameter-free 3D spatial-energy attention into all nine residual blocks (Block-SimAM) achieves the peak performance of $71.50\%$ Moderate mAP and $73.45\%$ $\text{mAP}_9$ ($+0.92\,\text{pp}$ over $A_7$) with negligible latency overhead.

While Table~\ref{tab:components} outlines this progressive single-factor roadmap, rigorous controlled micro-ablations isolating the individual contributions of loss objectives and attention under identical experimental budgets are analyzed below.


\paragraph{Geometric Supervision and Quality Head Factorial Study.}
Table~\ref{tab:loss} presents a standardized $2\times2$ factorial study evaluating geometric supervision (Smooth-L1 vs. MGIoU vs. OGA), quality prediction (with and without an explicit IQA head), and the head-free Q-OGA formulation under a strictly controlled 50-epoch schedule on NVIDIA L4. All configurations share the identical standardized depth $[3, 4, 2]$ backbone, Strata-BEV14 input, SG-FPN neck, decoupled grouped heads, and Block-SimAM attention ($0.754\,\text{M}$ parameters with IQA, $0.735\,\text{M}$ without IQA). Figure~\ref{fig:loss_ablation_curves} displays the training loss, validation loss, and Moderate mAP trajectories across the 50-epoch training schedule.

\begin{table*}[t]
\centering\small
\caption{Ablation of geometric supervision (Baseline Smooth-L1 vs. MGIoU vs. OGA), localization quality prediction (with and without an explicit IQA Head), and head-free Q-OGA on KITTI validation at Moderate difficulty ($R_{40}$) under a strictly standardized 50-epoch schedule on NVIDIA L4. All configurations share the standardized depth $[3, 4, 2]$ backbone with parameter-free Block-SimAM, Strata-BEV14 input, SG-FPN neck, and decoupled grouped heads. IQA models rank NMS candidates using composite score $S_{\mathrm{det}} = \sqrt{P_{\mathrm{cls}} \hat{q}}$. Best 50-epoch values are in \textbf{bold}, and second-best are \underline{underlined}.}
\label{tab:loss}
\setlength{\tabcolsep}{4pt}
\begin{tabular*}{\textwidth}{@{\extracolsep{\fill}}>{\raggedright\arraybackslash}p{0.32\textwidth}ccrrrr@{}}
\toprule
\textbf{Training Objective} & \textbf{\shortstack{IQA\\Head}} & \textbf{\shortstack{Params\\(M) $\downarrow$}} & \textbf{\shortstack{Car AP$_{\mathrm{M}}$\\(\%) $\uparrow$}} & \textbf{\shortstack{Pedestrian\\AP$_{\mathrm{M}}$ (\%) $\uparrow$}} & \textbf{\shortstack{Cyclist\\AP$_{\mathrm{M}}$ (\%) $\uparrow$}} & \textbf{\shortstack{mAP$_{\mathrm{M}}$\\(\%) $\uparrow$}} \\
\midrule
Baseline Smooth-L1 & No  & \textbf{0.735} & 81.09 & 54.37 & 64.99 & 66.82 \\
MGIoU Only (No Corner Loss) & No  & \textbf{0.735} & \textbf{84.91} & 54.00 & 64.18 & 67.69 \\
OGA (MGIoU + $\pi$-Corner) & No  & \textbf{0.735} & 83.91 & 58.38 & 65.62 & 69.30 \\
Baseline Smooth-L1 & Yes & 0.754 & 83.01 & \textbf{61.15} & 65.72 & \textbf{69.96} \\
\textbf{OGA (MGIoU + $\pi$-Corner)} (50 ep) & \textbf{Yes} & 0.754 & \underline{84.68} & 57.56 & \underline{66.26} & \underline{69.50} \\
\midrule
Q-OGA (Quality-Coupled Focal Loss) & No & \textbf{0.735} & 84.19 & 57.37 & 65.72 & 69.09 \\
\midrule
\multicolumn{7}{@{}l}{\textit{Standard Convergence Schedule (100 epochs)}} \\
\textbf{MobilePIXORNeXt (Final Checkpoint)} & \textbf{Yes} & 0.754 & 86.90 & 60.75 & 66.73 & 71.46 \\
\bottomrule
\end{tabular*}
\par\vspace{2pt}
{\footnotesize \textit{Note.} All 50-epoch variants are trained under identical conditions on NVIDIA L4 using standardized depth $(3, 4, 2)$, Strata-BEV14 representation, SG-FPN neck, decoupled grouped heads, and Block-SimAM attention ($0.754\,\text{M}$ parameters with IQA, $0.735\,\text{M}$ without IQA). MGIoU-only increases Car AP by $+3.82\,\mathrm{pp}$ over Smooth-L1, while adding $\pi$-Corner boosts Pedestrian AP by $+4.38\,\mathrm{pp}$, demonstrating complementary geometric benefits. Q-OGA integrates quality directly into classification supervision, matching explicit quality performance while saving $18{,}626$ parameters ($18.63\,\text{K}$). The final reported model completes 100 epochs to reach $71.46\%$ ($71.50\%$ rounded) Moderate mAP.}
\end{table*}


\begin{figure*}[t]
\centering
\paperfigure[width=\textwidth]{loss_ablation_curves.pdf}{42mm}{Multi-task training and validation convergence trajectories for geometric supervision and quality prediction ablation variants in Table~\ref{tab:loss}.}
\captionof{figure}{Multi-task optimization dynamics and validation accuracy trajectories across the 50-epoch training horizon for the geometric supervision and quality estimation variants in Table~\ref{tab:loss}. \textbf{(a)} Total training objective descent; \textbf{(b)} Total validation loss convergence; \textbf{(c)} Validation Moderate mAP ($R_{40}$, \%) trajectories evaluated at scheduled validation checkpoints. Oriented Geometric Alignment (OGA) and Q-OGA demonstrate superior optimization stability and accelerated convergence toward the peak performance regime.}
\label{fig:loss_ablation_curves}
\end{figure*}

Without IQA, replacing baseline Smooth-L1 ($66.82\%$ Moderate mAP) with full OGA ($69.30\%$) yields a substantial $+2.48\,\text{pp}$ improvement. Dissecting OGA into its constituent terms reveals complementary geometric specialization across categories:
\begin{itemize}
    \item \textbf{MGIoU Only}: Optimizing bounding box overlap via multi-axis projection GIoU alone raises Moderate mAP to $67.69\%$, delivering a dramatic $+3.82\,\text{pp}$ boost on rigid vehicles (Car AP$_{\mathrm{M}}$ increases from $81.09\%$ to $84.91\%$), while pedestrian accuracy remains virtually flat ($54.37\% \to 54.00\%$).
    \item \textbf{$\pi$-Corner Supervision}: Augmenting MGIoU with $\pi$-symmetric corner loss in full OGA elevates Pedestrian AP$_{\mathrm{M}}$ from $54.00\%$ to $58.38\%$ ($+4.38\,\text{pp}$) and Cyclist AP$_{\mathrm{M}}$ from $64.18\%$ to $65.62\%$ ($+1.44\,\text{pp}$), demonstrating that direct point-to-point boundary alignment resolves orientation ambiguity on small, deformable vulnerable road users.
\end{itemize}

When explicit localization quality estimation is introduced, Baseline Smooth-L1 reaches $69.96\%$ Moderate mAP, while OGA reaches $69.50\%$ at 50 epochs. Under the full 100-epoch schedule, MobilePIXORNeXt (OGA with IQA and Block-SimAM) achieves the peak performance of $71.46\%$ ($71.50\%$ rounded) Moderate mAP ($86.90\%$ Car, $60.75\%$ Pedestrian, $66.73\%$ Cyclist) and $73.45\%$ $\text{mAP}_9$.

Furthermore, Table~\ref{tab:loss} evaluates \textbf{Q-OGA}, which directly couples predicted classification confidence with the detached geometric quality target during training, eliminating the dedicated IQA prediction head. Q-OGA attains $69.09\%$ Moderate mAP, trailing the 50-epoch explicit IQA model by only $0.41\,\text{pp}$ while saving $18{,}626$ parameters ($18.63\,\text{K}$) and reducing deployable parameters to $734{,}943$ ($0.735\,\text{M}$). This validates Q-OGA as an ultra-compact deployment option for edge accelerators with constrained memory budgets.

\paragraph{3D Spatial-Energy Attention Integration and Placement.}
Table~\ref{tab:attention} evaluates the placement and topological integration of parameter-free SimAM 3D spatial-energy attention across the network hierarchy on NVIDIA L4.

\begin{table*}[t]
\centering\normalsize
\caption{Ablation of 3D spatial-energy attention (SimAM) integration topologies and placement strategies within the MobilePIXORNeXt backbone on the KITTI validation benchmark ($R_{40}$) on NVIDIA L4. Stage-level variants append a single attention adapter at stage downsampling boundaries ($C_2, C_3, C_4$), whereas the proposed Block-SimAM embeds parameter-free 3D energy attention inside all nine \texttt{MobilePIXORNeXt-S} residual blocks before LayerScale scaling. The no-attention baseline is strictly matched to the final depth $[3, 4, 2]$ and 100-epoch schedule. Best values are in \textbf{bold}, and second-best are \underline{underlined}.}
\label{tab:attention}
\setlength{\tabcolsep}{4pt}
\begin{tabular}{@{}llrrrr@{}}
\toprule
\textbf{\shortstack[l]{SimAM Attention\\Placement}} & \textbf{\shortstack[l]{Neck Fusion\\Width}} & \textbf{\shortstack{mAP$_{\mathrm{M}}$\\(\%) $\uparrow$}} & \textbf{\shortstack{mAP$_{9}$\\(\%) $\uparrow$}} & \textbf{\shortstack{Model\\Stage (ms) $\downarrow$}} & \textbf{\shortstack{Model\\FPS $\uparrow$}} \\
\midrule
None (No attention, 100 ep) & $F_{32}$ & 68.94 & 71.35 & \underline{6.82} & \underline{146.60} \\
Stage-2 ($C_2$) & $F_{32}$ & 69.44 & 71.49 & 7.30 & 136.91 \\
Stage-3 ($C_3$) & $F_{24}$ & \underline{69.80} & \underline{71.98} & \textbf{6.94} & \textbf{144.05} \\
Stage-4 ($C_4$) & $F_{32}$ & 69.36 & 71.34 & 7.44 & 134.43 \\
\textbf{Block-level ($C_2, C_3, C_4$)} & \textbf{$F_{32}$} & \textbf{71.50} & \textbf{73.45} & 8.67 & 115.32 \\
\bottomrule
\end{tabular}
\par\vspace{2pt}
{\footnotesize \textit{Note.} The no-attention configuration uses the standardized depth $(3, 4, 2)$ backbone, SG-FPN neck, decoupled grouped heads, OGA+IQA, and 100-epoch schedule on NVIDIA L4, cleanly isolating the $+2.52\,\mathrm{pp}$ Moderate mAP ($+2.07\,\mathrm{pp}$ $\mathrm{mAP}_9$) gain of parameter-free Block-SimAM without depth or schedule confounding. The 3D energy computation incurs only $1.85\,\mathrm{ms}$ model-stage overhead ($6.82\,\mathrm{ms}$ vs. $8.67\,\mathrm{ms}$).}
\end{table*}

To establish a rigorous, uncontaminated baseline, the no-attention configuration is benchmarked using the identical standardized depth $[3, 4, 2]$ backbone, SG-FPN neck, decoupled grouped heads, and 100-epoch schedule. Under this matched protocol, the no-attention baseline achieves $68.94\%$ Moderate mAP and $71.35\%$ $\mathrm{mAP}_9$ with a model-stage latency of $6.82\,\text{ms}$ ($146.60\,\text{FPS}$). 

Embedding parameter-free Block-SimAM across all nine residual blocks boosts performance to $71.46\%$ ($71.50\%$ rounded) Moderate mAP and $73.42\%$ ($73.45\%$ rounded) $\mathrm{mAP}_9$, representing an isolated net gain of $+2.52\,\text{pp}$ Moderate mAP ($+2.07\,\text{pp}$ $\mathrm{mAP}_9$). Because SimAM derives 3D energy weights analytically from local neuron statistics without learnable parameters, this performance surge requires zero additional memory for weights. The analytical energy computation incurs a modest model-stage overhead of only $1.85\,\text{ms}$ ($6.82\,\text{ms} \to 8.67\,\text{ms}$).

In contrast, inserting single attention adapters at stage transition boundaries ($C_2, C_3, C_4$) yields inferior gains: Stage-2 ($69.44\%$), Stage-3 ($69.80\%$), and Stage-4 ($69.36\%$). This indicates that applying energy modulation only at downsampled stage transitions is insufficient to preserve fine-grained vertical and spatial features. Layer-by-layer modulation inside every residual block before LayerScale scaling is essential for suppressing background noise and sharpening foreground target saliency across all three classes.

\paragraph{Multi-Task Optimization Dynamics and Convergence.}
Figure~\ref{fig:loss_curves} illustrates the training and validation loss trajectories across the first 50 epochs for the progressive architectural ablation steps $A_0\text{--}A_8$ established in Table~\ref{tab:components}. All individual loss components—including classification, coordinate offset, bounding box dimension, doubled yaw angle, and geometric alignment—exhibit smooth, stable descent throughout optimization. The Temperature-Softmax Bounded Uncertainty Weighting (T-SBUW) mechanism effectively balances these diverse task losses on a conserved simplex, preventing numerical divergence or dominance by large-magnitude regression gradients without requiring manual hyperparameter retuning.

\paragraph{Convergence Validation and Checkpoint Selection.}
During training across 100 epochs, MobilePIXORNeXt exhibits rapid early convergence and stabilizes in the $71.0\text{--}71.5\%$ Moderate mAP range. Comparing the AP-selected checkpoint at epoch 100 ($71.50\%$ Moderate mAP, $73.45\%$ $\mathrm{mAP}_9$) with the loss-selected checkpoint at epoch 47 ($71.53\%$ mAP, $73.61\%$ $\mathrm{mAP}_9$) confirms consistent asymptotic accuracy. The epoch-100 model is selected for primary deployment due to superior rigid-vehicle localization ($86.90\%$ vs. $82.70\%$ Car Moderate AP) and tighter detection candidate clustering, which reduces post-processing NMS latency by $32.99\%$ ($37.29\,\text{ms}$ vs. $55.65\,\text{ms}$).


\begin{figure*}[p]
\centering
\paperfigure[width=\textwidth,height=1.0\textheight,keepaspectratio]{training_curves_all_losses.pdf}{0.88\textheight}{Training and validation loss trajectories across 50 epochs.}
\captionof{figure}{Training and validation loss trajectories across the first 50 epochs under multi-task optimization for progressive architectural milestones ($A_0\text{--}A_8$). Panels display overall loss descent alongside individual terms: classification, geometric alignment, localization quality, coordinate offset, bounding box dimension, and doubled yaw angle. Losses exhibit stable convergence without gradient explosion or task starvation under T-SBUW.}
\label{fig:loss_curves}
\end{figure*}



\subsection{Qualitative Visualizations and Analysis}
\label{sec:qualitative_analysis}

\begin{figure*}[!t]
\centering
\paperfigure[width=\textwidth]{comparison_mobilepixor_heatmap.pdf}{65mm}{Comparative BEV detection and classification heatmap visualizations between MobilePIXOR and MobilePIXORNeXt.}
\captionof{figure}{Qualitative comparison of Bird's-Eye View (BEV) detections and predicted classification heatmaps between the baseline MobilePIXOR and MobilePIXORNeXt (Ours) across two representative driving scenes. \textbf{Row A}: MobilePIXOR exhibits diffused, low-amplitude heatmap responses with noticeable background clutter, leading to missed detections (false negatives) on distant and sparse point clusters. \textbf{Row B}: MobilePIXORNeXt produces sharp, concentrated focal peaks strictly centered on vehicle targets with suppressed background activations, successfully recovering distant objects and yielding tightly aligned oriented bounding footprints.}
\label{fig:comparison_heatmap}
\end{figure*}

\paragraph{Comparative Heatmap and Detection Analysis against Baseline.}
Figure~\ref{fig:comparison_heatmap} provides a side-by-side qualitative comparison between the baseline MobilePIXOR and MobilePIXORNeXt across two challenging driving scenarios, illustrating both oriented BEV bounding footprints and the corresponding foreground classification heatmaps.

Two key observations emerge from this qualitative comparison:
\begin{itemize}
    \item \textbf{Sharp Foreground Heatmap Concentration}: MobilePIXORNeXt exhibits sharp, concentrated confidence peaks tightly centered on vehicle targets with minimal background clutter. In contrast, baseline MobilePIXOR produces diffused activations with notable false-positive responses in background vegetation and road boundaries.
    \item \textbf{Distant and Sparse Target Recovery}: On distant targets characterized by sparse LiDAR reflection points, MobilePIXORNeXt successfully recovers valid detection peaks and tightly oriented bounding footprints, whereas MobilePIXOR suffers from missed detections (false negatives).
\end{itemize}

\paragraph{Visual Inspection on Representative Driving Scenarios.}
Figure~\ref{fig:qualitative_detections} presents MobilePIXORNeXt predictions on three KITTI validation frames (top to bottom: \texttt{007038}, \texttt{005571}, and \texttt{006761}). Class labels and confidence scores are retained in the camera overlays to support inspection of the predicted footprints; the BEV panels retain class annotations.

\begin{figure*}[!t]
\centering
\paperfigure[width=2.0\textwidth,height=0.4\textheight,keepaspectratio]{qualitative_detections.pdf}{85mm}{Qualitative BEV detection visualizations on KITTI validation scenes.}
\captionof{figure}{Qualitative BEV predictions from MobilePIXORNeXt (converged MobilePIXORNeXt deployment model) on KITTI validation frames \texttt{007038}, \texttt{005571}, and \texttt{006761} (top to bottom). \textbf{Left}: Front-camera images with calibrated projections of predicted BEV footprints, class labels, and confidence scores. \textbf{Right}: Corresponding LiDAR BEV views with predicted oriented footprints and class annotations. Prediction colors are cyan for Car, magenta for Pedestrian, and yellow for Cyclist. The panels visualize predicted footprints rather than full 3D boxes.}
\label{fig:qualitative_detections}
\end{figure*}

The paired camera and BEV views in Figure~\ref{fig:qualitative_detections} support inspection of the following prediction patterns:
\begin{itemize}
    \item \textbf{Vehicle Predictions}: Frame \texttt{007038} contains several predicted vehicle footprints at different forward distances, visible in both the camera projection and the BEV panel.
    \item \textbf{Mixed-Class Predictions}: Frame \texttt{005571} contains Car, Pedestrian, and Cyclist predictions. The retained labels and confidence scores identify the classes assigned to nearby footprints.
    \item \textbf{Nearby Road Users}: Frame \texttt{006761} shows predicted footprints around vehicles, pedestrians, and cyclists in the same scene. The two views reveal the spatial arrangement of these predictions.
\end{itemize}

\subsection{Discussion and Practical Limitations}
\label{sec:discussion_limitations}

\paragraph{Discussion and Insights.}
Our empirical evaluation demonstrates that high-precision oriented BEV footprint detection can be achieved with fewer than $0.76\,\text{M}$ deployable parameters without relying on complex PointNet encoders or memory-intensive 3D sparse convolutions. Deterministic statistical encoding (Strata-BEV14) effectively retains continuous height distribution and subcell centroid information on a regular 2D grid. The combination of grouped prediction heads, physics-consistent data augmentation (PC-GTA), and oriented geometric alignment (OGA) enables the lightweight ConvNeXt backbone to achieve competitive accuracy against significantly larger anchor-based 3D detectors while maintaining a sub-10\,ms model-stage inference time on modern GPU hardware.

\paragraph{Limitations and Future Scope.}
Several practical constraints offer directions for future research. First, MobilePIXORNeXt focuses on horizontal bounding footprints and heading angles modulo $\pi$, omitting full 3D vertical extents and directed orientation vectors. Second, the current benchmark is evaluated on 64-beam mechanical LiDAR data under clear weather conditions; evaluating robustness against sparse low-beam sensors (e.g., 16-beam or solid-state LiDAR), non-planar terrains, and adverse weather conditions remains an important next step. Finally, while model-stage execution is exceptionally fast ($8.67\,\text{ms}$, $115.32\,\text{FPS}$), end-to-end latency is partially bounded by CPU rasterization and sequential host-to-device memory transfer. Future work will investigate custom CUDA kernel parallelization for Strata-BEV14 encoding and int8/fp8 quantization on embedded edge processors such as NVIDIA Jetson Orin.
